%! Linear Combinations of Vectors — Unit 1, Lesson 1.1 — Exit Ticket (Answer Key)
\documentclass[10pt]{article}
\usepackage{linalg-article}
\usepackage{linalg-key}   % pulls in -boxes; do NOT also load -boxes
\newcommand{\TallMath}[1]{$\displaystyle #1\rule[-1.4em]{0pt}{3.2em}$}
\newcommand{\vv}{\mathbf{v}}
\newcommand{\ww}{\mathbf{w}}

\begin{document}

\pageheader{Unit 1, Lesson 1.1}{Exit Ticket --- Answer Key}
\namedateperiod

\vspace{0.15in}

Let $\vv = \begin{bmatrix} 3 \\ 1 \end{bmatrix}$ and $\ww = \begin{bmatrix} 1 \\ 2 \end{bmatrix}$.
Show your work.

\begin{enumerate}[leftmargin=1.8em, itemsep=14pt]
  \item Compute the linear combination $2\vv + 3\ww$.
        \ansline{$2\begin{bsmallmatrix} 3\\1 \end{bsmallmatrix}+3\begin{bsmallmatrix} 1\\2 \end{bsmallmatrix}
        = \begin{bsmallmatrix} 6\\2 \end{bsmallmatrix}+\begin{bsmallmatrix} 3\\6 \end{bsmallmatrix}
        = \begin{bsmallmatrix} 9\\8 \end{bsmallmatrix}$}
  \item Find weights $c$ and $d$ so that $c\,\vv + d\,\ww = \begin{bmatrix} 5 \\ 5 \end{bmatrix}$.
        \ansline{$3c+d=5,\ c+2d=5 \Rightarrow c=1,\ d=2$; check $\vv+2\ww=(3,1)+(2,4)=(5,5)$}
  \item Sam says the two vectors $\begin{bmatrix} 2 \\ 4 \end{bmatrix}$ and
        $\begin{bmatrix} 1 \\ 2 \end{bmatrix}$ can reach \emph{every} point in the plane.
        Is Sam right? Explain.
        \ansline{No --- $\begin{bsmallmatrix} 2\\4 \end{bsmallmatrix}=2\begin{bsmallmatrix} 1\\2 \end{bsmallmatrix}$,
        so the two are parallel. Every combination lands on the single line $y=2x$, so most
        points in the plane are unreachable.}
\end{enumerate}

\begin{teachernote}
Item 3 checks the central misconception --- assuming \emph{any} two vectors span the plane.
Full credit needs both ``they are parallel / one is a multiple of the other'' \emph{and} the
consequence ``they span only a line.''
\end{teachernote}

\end{document}
